\renewcommand{\baselinestretch}{1.5}\normalsize
\chapter{Introduction}

\section{Context and Operational Challenges}

The paradigm of higher education assessment has transitioned substantially toward digital and distributed examination delivery. While remote testing infrastructure expands accessibility and reduces institutional scheduling overhead, it fundamentally dismantles the environmental assumptions of classical physical examination halls: invariant seating arrangements, direct line-of-sight visual invigilation, acoustic isolation, and physical perimeter controls over unauthorised reference media.

In an unconstrained remote environment, academic dishonesty manifests across diverse, concurrent channels:
\begin{itemize}
    \item \textbf{Visual and Physical Evasion}: Candidacy substitution, co-presence of unauthorized secondary personnel within or adjacent to the camera frustum, placement of mobile communication devices outside standard webcam viewing angles, and physical concealment of course literature.
    \item \textbf{Acoustic Collaboration}: Whispered verbal exchanges, audio prompt broadcasting from external accomplices, or playback from synthesized speech tools.
    \item \textbf{Peripheral and OS Interaction}: Rapid context-switching between the examination interface and external browser windows, automated code generation tools, or direct clipboard injection of externally synthesized solutions.
    \item \textbf{Hardware and Stream Tampering}: Intentional sensor occlusion, lens blinding, or video loop replay attacks designed to spoof standard frame-liveness checks.
\end{itemize}

Manual human invigilation via multi-stream video grids fails at scale; cognitive load studies demonstrate marked vigilance decrement when a single human supervisor monitors more than four to six concurrent streams for intervals exceeding 45 minutes \cite{vigilance1948, proctoru2020}. Conversely, legacy automated lockdown software relies on static heuristics that trigger excessive false-positive rates on innocent physical micro-behaviors (such as keyboard-directed gaze) while failing to detect sophisticated multimodal violations.

\section{Problem Statement}

We formally define the online proctoring challenge as a real-time, multi-modal anomaly detection and decision fusion problem:

\begin{quote}
    \textit{Given continuous, asynchronous streams of high-frequency visual telemetry ($\mathcal{V}$ at 15\,FPS), raw acoustic pulse-code modulation ($\mathcal{A}$ at 16\,kHz), and discrete operating system interaction events ($\mathcal{E}$), construct a distributed multi-agent system capable of generating a calibrated composite risk metric $r(t) \in [0, 1]$ with an end-to-end alert latency $\tau \leq 1000$\,ms, while simultaneously providing deterministic, human-interpretable natural-language evidentiary traces $\mathcal{T}(t)$ and immutable cryptographic integrity receipts $\mathcal{R}$.}
\end{quote}

The core technical complexities encompass:
\begin{enumerate}
    \item \textbf{Asynchronous Modality Synchronization}: Reconciling continuous frame buffers with sporadic OS blur and clipboard events without introducing queue bloat or backpressure stalls.
    \item \textbf{Noise-Robust Perception}: Discriminating innocent posture adjustments and environmental background noise from deliberate cheating maneuvers.
    \item \textbf{Evidentiary Traceability}: Ensuring automated flags are backed by inspectable rationale strings rather than opaque, uncalibrated scalar outputs.
\end{enumerate}

\section{Research Objectives}

To resolve these challenges, the project establishes the following specific research and engineering objectives:

\begin{enumerate}
    \item \textbf{Distributed Monorepo Pipeline (O1)}: Engineer an integrated multi-tier microservices architecture using Turborepo to decouple compute-intensive perception from real-time alert dispatching and user interface rendering.
    \item \textbf{High-Throughput Vision Guard (O2)}: Construct a dedicated perception service wrapping YOLOv11 and MediaPipe Face Mesh pipelines to perform continuous multi-face, device, and posture inference at sub-70\,ms per-frame compute budgets.
    \item \textbf{Continuous Acoustic Verification (O3)}: Implement voice activity detection algorithms capable of dynamically adapting to background acoustic noise floors in home testing environments.
    \item \textbf{Behavioral Biometric Correlation (O4)}: Develop client-side telemetry listeners to quantify clipboard volume anomalies via Z-score metrics and detect elevated OS focus-switch frequencies.
    \item \textbf{Dynamic Decision Orchestration (O5)}: Formulate a weighted linear aggregation framework with exponential moving average (EMA) smoothing ($\alpha = 0.3$) to eliminate high-frequency score jitter.
    \item \textbf{Explainable AI (XAI) Synthesis (O6)}: Implement deterministic rationale generation mapping raw sensor activations into structured, non-technical explanatory summaries for human invigilators.
    \item \textbf{Cryptographic Audit Ledger (O7)}: Implement post-session receipt generation anchoring submission payloads with SHA-256 hash chains to ensure non-repudiation and regulatory compliance.
\end{enumerate}

\section{Methodological Scope}

The scope of this investigation encompasses:
\begin{itemize}
    \item Client-side video, audio, and browser-event acquisition within Chromium-compliant web browsers via standard WebRTC and Web Audio APIs.
    \item Server-side pipeline processing across nine dedicated microservices (Ports 4000--4009) and three specialized web interfaces (Ports 3000--3003).
    \item Empirical validation against an annotated benchmark of 200 standardized evaluation sessions encompassing 4,800 labelled 30-second inspection windows.
\end{itemize}

Explicitly excluded from the current scope are hardware-level side-channel attacks (e.g., HDMI splitters), kernel-level OS virtualization bypasses, and natural language semantic generation analysis of final essay answers.

\section{Report Organization}

The structure of this report is organized as follows:
\begin{itemize}
    \item \textbf{Chapter 2 (Literature Review)}: Synthesizes prior academic literature in automated proctoring, deep visual detectors, gaze estimation, audio diarization, and multimodal decision fusion, contrasting existing limitations against SentinelAI.
    \item \textbf{Chapter 3 (System Design and Methodology)}: Details the distributed monorepo architecture, individual agent processing mechanics, mathematical fusion models, database schemas, and evaluation protocols.
    \item \textbf{Chapter 4 (System Workflow Flowchart)}: Details the complete end-to-end operational execution path from session initialization to cryptographic receipt storage.
    \item \textbf{Chapter 5 (Results and Analysis)}: Presents quantitative empirical findings, precision-recall trade-offs, ablation outcomes, latency benchmarks, and false-positive analyses.
    \item \textbf{Chapter 6 (Conclusion and Future Work)}: Summarizes research milestones, outlines architectural constraints, and identifies prospective enhancements in adaptive reinforcement learning fusion and federated deployment.
\end{itemize}
